%% toc: P2 16 16.5,16.6,16.7
\chapter{Eliminating a Defined Relation Symbol}
\label{sec:p02-w02b-p02-018-defined-relation-elimination}

A definitional extension adds a relation symbol \(R\) to a signature
\(\Lang\) together with an axiom \(\delta_R\) saying that \(R\) holds exactly
where a fixed \(\Lang\)-formula \(\varphi\) holds (Example ``Definitions
inside a theory'' in \cref{ch:fol-theories}).
\Cref{thm:definitional-conservativity} shows that every \(\Lang\)-model of
the theory has exactly one expansion satisfying \(\delta_R\), and that the
extended theory proves no new sentence of the old signature.  This chapter
supplies the remaining piece: a translation \(\psi\mapsto\psi^{*}\) that
removes \(R\) from \emph{every} formula of the extended signature, and a
proof that the unique expansion satisfies \(\psi\) exactly when the old
structure satisfies \(\psi^{*}\), under every assignment.

\emph{Inputs.}  A theory \(T\) in \(\Lang\), a new \(n\)-ary relation symbol
\(R\), pairwise distinct variables \(x_1,\ldots,x_n\), and an
\(\Lang\)-formula \(\varphi\) whose free variables lie among them.
\emph{Conclusions.}  For every formula \(\psi\) of the extended signature,
an \(\Lang\)-formula \(\psi^{*}\) with no more free variables than \(\psi\)
(\cref{thm:p02-w02b-p02-018-translation-syntax}); agreement of \(\psi\) in
the expansion with \(\psi^{*}\) in the original structure
(\cref{thm:p02-w02b-p02-018-satisfaction-translation}); and the consequence
and derivability forms \(\{\delta_R\}\entails\psi\leftrightarrow\psi^{*}\)
and \(T\cup\{\delta_R\}\entails\psi\) iff \(T\entails\psi^{*}\)
(\cref{thm:p02-w02b-p02-018-consequences}).

\section{Standing data and the defined relation}
\label{sec:p02-w02b-p02-018-standing}

\begin{hbconvention}{def:p02-w02b-p02-018-standing}{Standing data for a definitional extension}
Throughout this chapter:
\begin{enumerate}
 \item \(\Lang\) is a signature, \(T\) is a theory in \(\Lang\), and
       \(R\) is a relation symbol not in \(\Lang\) of arity \(n\geq0\);
 \item \(x_1,\ldots,x_n\) are pairwise distinct variables, and
       \(\varphi\in\mathsf{Fm}(\Lang)\) satisfies
       \(\FV(\varphi)\subseteq\{x_1,\ldots,x_n\}\);
 \item \(\Lang_R\) is the signature obtained from \(\Lang\) by adjoining
       \(R\) with arity \(n\), and
       \[
        \delta_R\defeq\forall x_1\cdots\forall x_n
        \bigl(R(x_1,\ldots,x_n)\leftrightarrow\varphi\bigr),
       \]
       where \(\leftrightarrow\) is the defined biconditional and, for
       \(n=0\), the atomic formula is the nullary symbol \(R\) and no
       quantifier is written.
\end{enumerate}
Since \(R\) is a relation symbol, \(\Lang\) and \(\Lang_R\) have the same
function symbols, so their terms coincide.  ``Least'' variable always
refers to the fixed enumeration of the variables used for fresh choices in
\cref{sec:fol-syntax-semantics-04}; the names \(x_1,\ldots,x_n\) here
denote the defining variables, which need not be the first members of that
enumeration.  A consequence \(\Gamma\entails\chi\) is evaluated over
structures of the signature to which \(\Gamma\cup\{\chi\}\) belongs:
\(\Lang_R\)-structures for \(T\cup\{\delta_R\}\), and \(\Lang\)-structures
for \(T\).  For an assignment \(s\) and a tuple
\(\bar a=(a_1,\ldots,a_n)\), write
\(s[\bar x\mapsto\bar a]\) for \(s[x_1\mapsto a_1]\cdots[x_n\mapsto a_n]\).
\end{hbconvention}

\begin{hbdefinition}{def:p02-w02b-p02-018-defined-relation}{Relation defined by \(\varphi\), and the defining expansion}
Let \(\mathcal M\) be an \(\Lang\)-structure.  The relation \emph{defined by
\(\varphi\) in \(\mathcal M\)} is
\[
 R^{\mathcal M}_{\varphi}\defeq
 \bigl\{\bar a\in(\dom\mathcal M)^n:
 \mathcal M,s[\bar x\mapsto\bar a]\entails\varphi
 \text{ for every assignment }s\text{ in }\mathcal M\bigr\}.
\]
The \emph{defining expansion} \(\mathcal M_{\varphi}\) is the
\(\Lang_R\)-structure with domain \(\dom\mathcal M\), the interpretations of
\(\mathcal M\) for the symbols of \(\Lang\), and
\(R^{\mathcal M_{\varphi}}\defeq R^{\mathcal M}_{\varphi}\).  Thus
\(\mathcal M_{\varphi}\mathbin{\upharpoonright}\Lang=\mathcal M\).  For
\(n=0\), \((\dom\mathcal M)^0\) has the empty tuple as its only member,
which lies in \(R^{\mathcal M}_{\varphi}\) iff the sentence \(\varphi\) is
true in \(\mathcal M\).
\end{hbdefinition}

By \cref{thm:p02-w02b-p02-018-instance}(3) below, the quantifier ``for
every assignment'' in this definition may be replaced by ``for some
assignment''; in the proof of \cref{thm:definitional-conservativity} the
same relation is written \(R^{*}\).
\Cref{thm:p02-w02b-p02-018-consequences}(1) shows that
\(\mathcal M_{\varphi}\) is the unique expansion of \(\mathcal M\)
satisfying \(\delta_R\), whether or not \(\mathcal M\) is a model of \(T\).

\section{Instances of the defining formula}
\label{sec:p02-w02b-p02-018-instances}

To eliminate an atomic formula \(R(t_1,\ldots,t_n)\), the arguments
\(t_1,\ldots,t_n\) have to be put in place of \(x_1,\ldots,x_n\) in
\(\varphi\) \emph{simultaneously}.  The handbook's substitution
\(\chi[t/x]\) replaces one variable at a time and avoids capture.  The
following procedure builds the simultaneous instance from it.

\begin{hbmethod}{def:p02-w02b-p02-018-instance}{The instance \(\varphi(t_1,\ldots,t_n)\)}
\emph{Input:} terms \(t_1,\ldots,t_n\in\mathsf{Term}(\Lang)\).
\begin{enumerate}
 \item \emph{Fresh names.}  Let \(z_1,\ldots,z_n\) be the \(n\) least
       variables, listed in increasing order, that lie outside
       \[
        V\defeq\operatorname{Var}(\varphi)\cup\{x_1,\ldots,x_n\}
        \cup\operatorname{Var}(t_1)\cup\cdots\cup\operatorname{Var}(t_n).
       \]
 \item \emph{Renaming stage.}  Put \(\chi_0\defeq\varphi\) and
       \(\chi_k\defeq\chi_{k-1}[z_k/x_k]\) for \(k=1,\ldots,n\).
 \item \emph{Insertion stage.}  Put \(\theta_0\defeq\chi_n\) and
       \(\theta_k\defeq\theta_{k-1}[t_k/z_k]\) for \(k=1,\ldots,n\).
 \item \emph{Output:} \(\varphi(t_1,\ldots,t_n)\defeq\theta_n\).
\end{enumerate}
\emph{Side conditions.}  The set \(V\) is finite and the variables are
infinitely many, so step~1 always succeeds; the choice is fixed by the
enumeration, so the output is a function of \(\varphi\), \(x_1,\ldots,x_n\)
and \(t_1,\ldots,t_n\).  The variables \(z_1,\ldots,z_n\) are pairwise
distinct, differ from every \(x_i\), and occur in no \(t_j\) and nowhere in
\(\varphi\).  Capture of a variable of some \(t_k\) by a quantifier of
\(\varphi\) is prevented by the renaming clause of capture-avoiding
substitution itself.  For \(n=0\) both stages are empty and
\(\varphi()\) is \(\varphi\).
\end{hbmethod}

\begin{hblemma}{thm:p02-w02b-p02-018-instance}{proved=annex:p02-w02b-p02-018-instance}{Instance lemma}
Let \(t_1,\ldots,t_n\in\mathsf{Term}(\Lang)\).
\begin{enumerate}
 \item \(\varphi(t_1,\ldots,t_n)\in\mathsf{Fm}(\Lang)\) and
       \(\FV(\varphi(t_1,\ldots,t_n))\subseteq
       \operatorname{Var}(t_1)\cup\cdots\cup\operatorname{Var}(t_n)\).
 \item For every structure \(\mathcal N\) whose signature extends
       \(\Lang\) and every assignment \(s\) in \(\mathcal N\), with
       \(a_i\defeq\sem{t_i}^{\mathcal N}_s\),
       \[
        \mathcal N,s\entails\varphi(t_1,\ldots,t_n)
        \quad\Longleftrightarrow\quad
        \mathcal N,s[\bar x\mapsto\bar a]\entails\varphi .
       \]
 \item For every \(\Lang\)-structure \(\mathcal M\), every assignment \(s\)
       in \(\mathcal M\) and every \(\bar a\in(\dom\mathcal M)^n\),
       \(\bar a\in R^{\mathcal M}_{\varphi}\) iff
       \(\mathcal M,s[\bar x\mapsto\bar a]\entails\varphi\).
\end{enumerate}
\end{hblemma}

\section{The translation and its satisfaction theorem}
\label{sec:p02-w02b-p02-018-translation}

\begin{hbdefinition}{def:p02-w02b-p02-018-translation}{Elimination of \(R\)}
For \(\psi\in\mathsf{Fm}(\Lang_R)\), the formula \(\psi^{*}\) is
defined by recursion on the complexity \(c(\psi)\):
\[
\begin{aligned}
 (r=t)^{*}&\defeq r=t,\\
 (S(t_1,\ldots,t_m))^{*}&\defeq S(t_1,\ldots,t_m)
   &&\text{for each relation symbol }S\in\Lang,\\
 (R(t_1,\ldots,t_n))^{*}&\defeq\varphi(t_1,\ldots,t_n),\\
 \bottom^{*}&\defeq\bottom,\qquad
 (\neg\chi)^{*}\defeq\neg\chi^{*},\\
 (\chi\circ\eta)^{*}&\defeq(\chi^{*}\circ\eta^{*})
   &&\text{for }\circ\in\{\land,\lor,\to\},\\
 (Qy\,\chi)^{*}&\defeq Qy\,\chi^{*}
   &&\text{for }Q\in\{\forall,\exists\}.
\end{aligned}
\]
The instance \(\varphi(t_1,\ldots,t_n)\) is that of
\cref{def:p02-w02b-p02-018-instance}.  By unique readability
(\cref{thm:fol-unique-readability}) exactly one clause applies to each
formula, and every recursive call is on an immediate subformula of lower
complexity, so \(\psi^{*}\) is well defined.  Abbreviations are expanded
before translating; in particular
\((\chi\leftrightarrow\eta)^{*}=(\chi^{*}\leftrightarrow\eta^{*})\).
\end{hbdefinition}

The translation changes only the atomic formulas whose relation symbol is
\(R\), and it leaves every quantifier and bound variable of \(\psi\) in
place.  The bound variables that capture-avoiding substitution introduces
are confined to the inserted instances of \(\varphi\).

\begin{hblemma}{thm:p02-w02b-p02-018-translation-syntax}{proved=annex:p02-w02b-p02-018-translation-syntax}{The translation lands in the old signature}
For every \(\psi\in\mathsf{Fm}(\Lang_R)\):
\begin{enumerate}
 \item \(\psi^{*}\in\mathsf{Fm}(\Lang)\) and
       \(\FV(\psi^{*})\subseteq\FV(\psi)\); in particular the translation
       of an \(\Lang_R\)-sentence is an \(\Lang\)-sentence;
 \item if \(\psi\in\mathsf{Fm}(\Lang)\), then \(\psi^{*}=\psi\).
\end{enumerate}
\end{hblemma}

\begin{hbtheorem}{thm:p02-w02b-p02-018-satisfaction-translation}{proved=annex:p02-w02b-p02-018-satisfaction-translation}{Satisfaction translation}
For every \(\Lang\)-structure \(\mathcal M\), every assignment \(s\) in
\(\mathcal M\) and every \(\psi\in\mathsf{Fm}(\Lang_R)\),
\[
 \mathcal M_{\varphi},s\entails\psi
 \quad\Longleftrightarrow\quad
 \mathcal M,s\entails\psi^{*}.
\]
In particular, if \(\mathcal M\) is an \(\Lang\)-model of \(T\), then
\(\mathcal M_{\varphi}\) is the unique \(\Lang_R\)-expansion of
\(\mathcal M\) satisfying \(\delta_R\), it is a model of
\(T\cup\{\delta_R\}\), and for every assignment \(s\) and every
\(\psi\in\mathsf{Fm}(\Lang_R)\) it satisfies \(\psi\) under \(s\)
iff \(\mathcal M\) satisfies \(\psi^{*}\) under \(s\).
\end{hbtheorem}

\begin{hbcorollary}{thm:p02-w02b-p02-018-consequences}{proved=annex:p02-w02b-p02-018-consequences}{Consequence and derivability forms}
\begin{enumerate}
 \item An \(\Lang_R\)-structure \(\mathcal N\) satisfies \(\delta_R\) iff
       \(R^{\mathcal N}=R^{\mathcal N\mathbin{\upharpoonright}\Lang}_{\varphi}\),
       that is, iff
       \(\mathcal N=(\mathcal N\mathbin{\upharpoonright}\Lang)_{\varphi}\).
       Hence every \(\Lang\)-structure has exactly one \(\Lang_R\)-expansion
       satisfying \(\delta_R\).
 \item For every \(\psi\in\mathsf{Fm}(\Lang_R)\),
       \(\{\delta_R\}\entails\psi\leftrightarrow\psi^{*}\), and therefore
       \(T\cup\{\delta_R\}\entails\psi\leftrightarrow\psi^{*}\).
 \item For every \(\Gamma\cup\{\psi\}\subseteq\mathsf{Fm}(\Lang_R)\), with
       \(\Gamma^{*}\defeq\{\gamma^{*}:\gamma\in\Gamma\}\),
       \[
        T\cup\{\delta_R\}\cup\Gamma\entails\psi
        \quad\Longleftrightarrow\quad
        T\cup\Gamma^{*}\entails\psi^{*}.
       \]
       In particular \(T\cup\{\delta_R\}\entails\psi\) iff
       \(T\entails\psi^{*}\).  For \(\psi\in\mathsf{Sent}(\Lang)\) and
       \(\Gamma=\varnothing\) this is clause~2 of
       \cref{thm:definitional-conservativity}, since then
       \(\psi^{*}=\psi\).
 \item Assume the Axiom of Choice.  For every
       \(\Gamma\cup\{\psi\}\subseteq\mathsf{Fm}(\Lang_R)\),
       \(T\cup\{\delta_R\}\derives_{\mathsf{ND}_{\mathsf{FOL}}}
       \psi\leftrightarrow\psi^{*}\), and
       \[
        T\cup\{\delta_R\}\cup\Gamma\derives_{\mathsf{ND}_{\mathsf{FOL}}}\psi
        \quad\Longleftrightarrow\quad
        T\cup\Gamma^{*}\derives_{\mathsf{ND}_{\mathsf{FOL}}}\psi^{*},
       \]
       the derivations on the left being in the signature \(\Lang_R\) and
       those on the right in \(\Lang\).
\end{enumerate}
\end{hbcorollary}

Clause~1 is the argument for clause~1 of
\cref{thm:definitional-conservativity}, now stated for every
\(\Lang\)-structure; clause~3 reduces any question about consequence in the
extended theory, from any further premises, to a question in the original
signature.  Clause~4 uses soundness (\cref{thm:fol-soundness}) and the
Axiom-of-Choice form of completeness (\cref{thm:fol-completeness}) in both
directions.

Kleene calls added symbols with their postulates \emph{eliminable} when an
effective process assigns to every formula of the enlarged system a formula
of the original one so that three relations hold: formulas of the original
system are left unchanged, each formula is provably equivalent in the
enlarged system to the formula assigned to it, and deducibility from
premises in the enlarged system passes to deducibility of the assigned
formulas from the assigned premises in the original system
\cite[\S\,74, p.~405]{Kleene1952}.  Taking the enlarged system to have the
premises \(T\cup\{\delta_R\}\) and the original one the premises \(T\),
\cref{thm:p02-w02b-p02-018-translation-syntax}(2) and clauses~2 and~4
above supply the counterparts of these three relations for the
natural-deduction calculus of this book under the Axiom of Choice.  The
translation is computed by a terminating recursion with the
deterministic fresh-variable choice of \cref{def:p02-w02b-p02-018-instance}.

\section{A worked translation}
\label{sec:p02-w02b-p02-018-example}

\begin{hbexample}{def:p02-w02b-p02-018-example}{Betweenness in a dense order}
Let \(\Lang=\{<\}\) with \(<\) binary, let \(T\) be the theory of dense
linear orders without endpoints (\cref{ch:fol-theories}), let \(R\) be
binary, and take distinct variables \(x_1,x_2,y\) with
\[
 \varphi\defeq\exists y\,(x_1<y\land y<x_2).
\]
In a model \(\mathcal M\) of \(T\), the relation
\(R^{\mathcal M}_{\varphi}\) holds of \((a,b)\) iff some element lies
strictly between \(a\) and \(b\); by density it coincides with
\(<^{\mathcal M}\).

\emph{The atom \(R(x_2,x_1)\).}  Here \(t_1=x_2\) and \(t_2=x_1\), so
\(V=\{x_1,x_2,y\}\) and \(z_1,z_2\) are the two least variables outside
\(V\).  The renaming stage gives
\(\chi_1=\exists y\,(z_1<y\land y<x_2)\) and
\(\chi_2=\exists y\,(z_1<y\land y<z_2)\); no binder is renamed because
\(y\) differs from the substituted variables and from the inserted ones.
The insertion stage gives \(\theta_1=\exists y\,(x_2<y\land y<z_2)\) and
\[
 (R(x_2,x_1))^{*}=\theta_2=\exists y\,(x_2<y\land y<x_1).
\]

\emph{The atom \(R(y,x_1)\), whose argument \(y\) is bound in
\(\varphi\).}  The renaming stage is as before.  In the insertion stage the
substituted term \(y\) contains the bound variable \(y\), so the
capture-avoiding clause renames that binder to the least variable \(w\)
outside \(\{z_1,z_2,y\}\) before inserting:
\(\theta_1=\exists w\,(y<w\land w<z_2)\).  The last step inserts \(x_1\)
for \(z_2\); if \(w\) happens to be \(x_1\), the binder is renamed once
more.  In either case the result has the shape
\[
 (R(y,x_1))^{*}=\exists w'\,(y<w'\land w'<x_1)
\]
for a bound variable \(w'\) distinct from \(y\) and \(x_1\) that the fixed
enumeration determines.  Consequently
\[
 \bigl(\forall x_1\,\exists y\,R(y,x_1)\bigr)^{*}
 =\forall x_1\,\exists y\,\exists w'\,(y<w'\land w'<x_1),
\]
and by \cref{thm:p02-w02b-p02-018-satisfaction-translation} the sentence
\(\forall x_1\,\exists y\,R(y,x_1)\) is true in \(\mathcal M_{\varphi}\)
iff every element of \(\mathcal M\) has two elements below it, which holds
in every model of \(T\) because there are no endpoints.  By
\cref{thm:p02-w02b-p02-018-consequences}(3),
\(T\cup\{\delta_R\}\entails\forall x_1\,\exists y\,R(y,x_1)\).
\end{hbexample}

\section{Failure cases and limits}
\label{sec:p02-w02b-p02-018-limits}

\begin{hbwarning}{def:p02-w02b-p02-018-naive-replacement}{Replacing the defining variables one at a time, or textually}
With \(\varphi=\exists y\,(x_1<y\land y<x_2)\) as in
\cref{def:p02-w02b-p02-018-example}, two shortcuts give wrong
translations.
\begin{enumerate}
 \item \emph{Sequential substitution without renaming.}  For
       \(R(x_2,x_1)\), computing \((\varphi[x_2/x_1])[x_1/x_2]\) first
       yields \(\exists y\,(x_2<y\land y<x_2)\) and then
       \(\exists y\,(x_1<y\land y<x_1)\).  This is false under every
       assignment in a strict order, whereas
       \((R(x_2,x_1))^{*}=\exists y\,(x_2<y\land y<x_1)\) is true in
       \((\mathbb Q,<)\) whenever \(x_2\) is assigned a number below the one
       assigned to \(x_1\).
       The second step altered the argument \(x_2\) inserted by the first;
       the renaming stage of \cref{def:p02-w02b-p02-018-instance} keeps the
       inserted arguments out of reach of later steps.
 \item \emph{Textual replacement.}  Writing \(y\) for \(x_1\) inside
       \(\varphi\) without renaming the binder gives
       \(\exists y\,(y<y\land y<x_2)\), in which the inserted \(y\) is
       captured; the formula is again false under every assignment in a
       strict order.  The capture-avoiding result computed in
       \cref{def:p02-w02b-p02-018-example} keeps the inserted \(y\) free.
\end{enumerate}
\end{hbwarning}

\begin{hbwarning}{def:p02-w02b-p02-018-limits}{What the translation does not cover}
\begin{enumerate}
 \item \emph{Without the defining axiom.}  The equivalence of \(\psi\) and
       \(\psi^{*}\) holds in \(\mathcal M_{\varphi}\) and, by
       \cref{thm:p02-w02b-p02-018-consequences}(1), in no other expansion.
       In the order \((\mathbb Q,<)\), the expansion interpreting \(R\) as
       the empty relation falsifies \(R(x_1,x_2)\) under an assignment with
       \(x_1\mapsto0\) and \(x_2\mapsto1\), while
       \((R(x_1,x_2))^{*}=\exists y\,(x_1<y\land y<x_2)\) is true there.
 \item \emph{Parameters in the defining formula.}  If \(\varphi\) had a
       free variable \(u\) outside \(\{x_1,\ldots,x_n\}\), the set of tuples
       satisfying \(\varphi\) would in general depend on the value of
       \(u\), so no single interpretation of \(R\) would serve every
       assignment, and \(\delta_R\) would not be a sentence.  The hypothesis
       \(\FV(\varphi)\subseteq\{x_1,\ldots,x_n\}\) is used in
       \cref{thm:p02-w02b-p02-018-instance}(3) and in
       \cref{thm:p02-w02b-p02-018-consequences}(1).
 \item \emph{Defined function and constant symbols.}  The chapter eliminates
       one relation symbol.  A function symbol occurs inside terms, possibly
       nested, so the atomic clause of
       \cref{def:p02-w02b-p02-018-translation} has no direct counterpart,
       and a defining formula for a function symbol must in addition be
       functional on the models of \(T\).  Nothing in this chapter proves a
       translation or conservativity result for that case.
\end{enumerate}
\end{hbwarning}

\section{Relationships and sources}
\label{sec:p02-w02b-p02-018-relationships}

The existence and uniqueness of the expansion and the conservativity of
\(T\cup\{\delta_R\}\) over \(T\) for \(\Lang\)-sentences are
\cref{thm:definitional-conservativity}; this chapter extends them to
every formula of \(\Lang_R\) through the translation, and recovers the
sentence case as \cref{thm:p02-w02b-p02-018-consequences}(3).  The proofs
use capture-avoiding substitution and its complexity and free-variable
lemmas (\cref{sec:fol-syntax-semantics-04}), the term-substitution lemma
(\cref{lem:term-substitution}), formula coincidence
(\cref{lem:fol-coincidence}), the substitution lemma
(\cref{thm:fol-substitution}) and invariance under reduct
(\cref{lem:reduct-invariance}).  The definitions, notation and conventions
are those of \cref{ch:fol-semantics} and \cref{ch:fol-theories}.  The
only external work cited is Kleene's definition of eliminability, used for
comparison after \cref{thm:p02-w02b-p02-018-consequences}; its page number
is that of the 1974 seventh reprint of \cite{Kleene1952}.  No source-specific
presentation of definitional extensions is reproduced in this chapter.
\index{definitional extension!translation}
\index{defined relation symbol!elimination}
